\documentclass{beamer}

\usepackage{ctex}
\usepackage{amsmath, amsfonts, amssymb, bm}
\usepackage{graphicx}
\usepackage{hyperref}
\usepackage{booktabs}
\usepackage{array}
\usepackage{ragged2e}

\usetheme{Berlin}
\usecolortheme{whale}
\useinnertheme[shadow=true]{rounded}
\usefonttheme[onlymath]{serif}
\setbeamercovered{transparent}
\setbeamertemplate{navigation symbols}{}

\setbeamertemplate{footline}{%
 \leavevmode%
 \hbox{%
 \begin{beamercolorbox}[wd=.31\paperwidth,ht=2.25ex,dp=1ex,center]{author in head/foot}
 \insertshortauthor
 \end{beamercolorbox}%
 \begin{beamercolorbox}[wd=.38\paperwidth,ht=2.25ex,dp=1ex,center]{title in head/foot}
 \insertshorttitle
 \end{beamercolorbox}%
 \begin{beamercolorbox}[wd=.31\paperwidth,ht=2.25ex,dp=1ex,center]{date in head/foot}
 \insertshortdate \quad \insertframenumber/\inserttotalframenumber
 \end{beamercolorbox}%
 }
}

\title[Matbench GNN 最终汇报]{晶体带隙预测实验最终汇报}
\author[刘嘉瑞]{汇报人：刘嘉瑞}
\institute[Band Gap]{项目：基于晶体图神经网络的 band gap 预测}
\date{2026年4月17日}

\begin{document}

\begin{frame}
 \titlepage
\end{frame}

\begin{frame}{项目目标}
\begin{itemize}
  \item 在鸿之微 \texttt{crystal\_band\_100000} 数据集上完成 6 组晶体 GNN 实验。
  \item 统一训练设定：单卡 RTX 3090、随机种子 42、所有实验强制训练 50 个 epoch。
  \item 任务定义：材料晶体 \textbf{band gap 回归预测}。
  \item 评价指标：MAE 为主，辅以 RMSE、$R^2$、Metal F1、训练时长和峰值显存。
\end{itemize}
\end{frame}

\begin{frame}{数据集与任务}
\small
\begin{itemize}
  \item 数据来源：鸿之微 \texttt{crystal\_band\_100000.jsonl}
  \item 清洗后有效样本：\textbf{95,497}
  \item 数据划分：
  \begin{itemize}
    \item 训练集：76,397
    \item 验证集：9,549
    \item 测试集：9,551
  \end{itemize}
  \item 主目标：\textbf{band\_gap}（eV）回归
  \item 可用辅助标签：\texttt{cbm}, \texttt{vbm}, \texttt{gap\_type}, \texttt{formula}, \texttt{density}
  \item E6 中进一步利用 CBM/VBM 构造辅助监督和物理一致性约束。
\end{itemize}
\end{frame}

\begin{frame}{实验设计}
\scriptsize
\centering
\setlength{\tabcolsep}{4pt}
\begin{tabular}{p{0.8cm} p{2.8cm} p{4.9cm} p{1.2cm}}
\toprule
编号 & 模型 & 核心改动 & 状态 \\
\midrule
E1 & CGCNN Baseline & 原子图卷积基线 & 完成 \\
E2 & ALIGNN Baseline & 原子图 + line graph，建模键角信息 & 完成 \\
E3 & M3GNet Baseline & 三体几何建模 & 完成 \\
E4 & ALIGNN + Virtual Node & 在 ALIGNN 上加入全局虚拟节点聚合与广播 & 完成 \\
E5 & ALIGNN + Transformer & 在 ALIGNN 表征后叠加 2 层 4 头 Transformer & 完成 \\
E6 & ALIGNN + Physics Constraint & 在 ALIGNN 上加入多头输出与物理一致性损失 & 完成 \\
\bottomrule
\end{tabular}

\vspace{2mm}
\footnotesize
说明：E4、E5、E6 是三个\textbf{分别独立}建立在 ALIGNN 基线上的增强实验，不是逐层叠加。
\end{frame}

\begin{frame}{最终结果总表}
\scriptsize
\centering
\resizebox{\textwidth}{!}{
\begin{tabular}{lcccccc}
\toprule
实验 & 模型 & MAE$\downarrow$ & RMSE$\downarrow$ & $R^2\uparrow$ & Metal F1$\uparrow$ & Best Epoch \\
\midrule
E1 & CGCNN & 0.302169 & 0.530287 & 0.884731 & 0.806929 & 48 \\
E2 & ALIGNN & 0.231432 & 0.471325 & 0.908939 & 0.862117 & 50 \\
E3 & M3GNet & 0.268041 & 0.553444 & 0.874444 & 0.871087 & 42 \\
E4 & ALIGNN + Virtual Node & \textbf{0.220077} & 0.474616 & 0.907663 & \textbf{0.883811} & 48 \\
E5 & ALIGNN + Transformer & 0.237450 & 0.491768 & 0.900869 & 0.876873 & 46 \\
E6 & ALIGNN + Physics Constraint & 0.221208 & \textbf{0.453924} & \textbf{0.915539} & 0.864048 & 50 \\
\bottomrule
\end{tabular}
}
\end{frame}

\begin{frame}{结果解读}
\small
\begin{itemize}
  \item \textbf{基线对比：}ALIGNN 明显优于 CGCNN 和 M3GNet，是当前最强基线。
  \item \textbf{按 MAE 排序：}E4 $>$ E6 $>$ E2 $>$ E5 $>$ E3 $>$ E1。
  \item \textbf{按 RMSE / $R^2$ 排序：}E6 最好，说明物理约束版整体误差分布更稳。
  \item \textbf{按 Metal F1 排序：}E4 最好，说明 Virtual Node 对金属/非金属判别更有帮助。
  \item \textbf{E5 现象：}Transformer 版没有超过 E4，也没有超过 E2，说明当前这版全局注意力设计尚未带来稳定收益。
\end{itemize}
\end{frame}

\begin{frame}{增强实验与 ALIGNN 基线对比}
\small
\begin{itemize}
  \item \textbf{E4 vs E2：}
  \begin{itemize}
    \item MAE 从 0.2314 降到 0.2201
    \item Metal F1 从 0.8621 提升到 0.8838
    \item 说明全局虚拟节点机制有效
  \end{itemize}
  \item \textbf{E5 vs E2：}
  \begin{itemize}
    \item MAE 从 0.2314 上升到 0.2375
    \item RMSE 也变差
    \item 说明当前 Transformer 版暂时没有带来收益
  \end{itemize}
  \item \textbf{E6 vs E2：}
  \begin{itemize}
    \item MAE 从 0.2314 降到 0.2212
    \item RMSE 从 0.4713 降到 0.4539
    \item $R^2$ 从 0.9089 提升到 0.9155
    \item 说明物理约束能明显提升整体回归稳定性
  \end{itemize}
\end{itemize}
\end{frame}

\begin{frame}{训练开销对比}
\scriptsize
\centering
\resizebox{\textwidth}{!}{
\begin{tabular}{lcccc}
\toprule
实验 & 训练轮数 & 训练时长（h） & 峰值显存（GB） & 备注 \\
\midrule
E1 CGCNN & 50 & 2.08 & 0.792 & 最轻量 \\
E2 ALIGNN & 50 & 5.74 & 12.329 & 最强基线 \\
E3 M3GNet & 50 & 3.16 & 3.147 & 三体建模 \\
E4 ALIGNN + VN & 50 & 6.18 & 12.210 & 最优 MAE / F1 \\
E5 ALIGNN + Transformer & 50 & 6.83 & 9.691 & 收益不明显 \\
E6 ALIGNN + Physics & 50 & 5.81 & 12.548 & 最优 RMSE / $R^2$ \\
\bottomrule
\end{tabular}
}

\vspace{2mm}
\small
\begin{itemize}
  \item 六组实验都完整跑满 50 个 epoch。
  \item ALIGNN 系列总体耗时更高，但精度明显优于 CGCNN。
  \item 数据缓存迁移到数据盘后，后续实验运行稳定，没有再次出现磁盘写满中断。
\end{itemize}
\end{frame}

\begin{frame}{E6 物理约束的含义}
\small
\[
L = L_{\text{bandgap}} + 0.3L_{\text{aux}} + 0.5L_{\text{phys}} + 0.1L_{\text{nonneg}}
\]
\begin{itemize}
  \item $L_{\text{bandgap}}$：主任务损失，让预测 band gap 尽量接近真实值。
  \item $L_{\text{aux}}$：辅助任务损失，让模型同时预测 CBM 和 VBM。
  \item $L_{\text{phys}}$：物理一致性损失，约束 \texttt{band gap = CBM - VBM}。
  \item $L_{\text{nonneg}}$：非负惩罚，避免预测出负带隙。
  \item 有效性 mask 规则：
  \[
  |cbm| > \varepsilon \;\; \text{or} \;\; |vbm| > \varepsilon
  \]
  仅在 CBM/VBM 明确有效的样本上计算辅助损失和物理约束，避免零占位样本污染训练。
\end{itemize}
\end{frame}

\begin{frame}{最终结论}
\begin{itemize}
  \item \textbf{最强基线：}ALIGNN
  \item \textbf{最优 MAE / Metal F1：}E4（ALIGNN + Virtual Node）
  \item \textbf{最优 RMSE / $R^2$：}E6（ALIGNN + Physics Constraint）
  \item \textbf{无效增强：}E5（ALIGNN + Transformer）在当前实现下没有优于 ALIGNN 基线
  \item 因此，本轮实验最重要的结论是：
  \begin{itemize}
    \item 全局虚拟节点增强是有效的
    \item 物理约束增强也是有效的
    \item Transformer 全局建模方案仍需继续调整
  \end{itemize}
\end{itemize}
\end{frame}

\begin{frame}{下一步工作}
\begin{itemize}
  \item 对 E5 做针对性调优，重点排查位置编码、全局 token 组织方式和注意力层深度。
  \item 进一步分析 E4 和 E6 的误差分布，找出它们各自提升最明显的材料类型。
  \item 若后续继续发论文或做答辩展示，可以以 E4 和 E6 作为两个主增强方向展开。
\end{itemize}
\end{frame}

\end{document}
